\subsection{Definitions of angles.}
In this section, we consider two definitions of the angle of a cloud.
In order to
give these definitions, we must first define the angle between spaces.
\\
In what follows, we shall assume that clouds have a nontrivial stabilizer and in
the notation for a cloud \( [M] \), \( M \) is its center.
\begin{definition}
    The \emph{angle} between spaces \( X_1, X_2 \in [M] \) such that \( |X_1
    M| = r_1 \), \( | X_2 M | =r_2 \), \( |X_1 X_2| = d\), where \( r_1,
    r_2 \neq 0 \), is the quantity \( \arccos \left(\frac{r_1^2 +
    r_{2}^2 - d^2}{2r_1r_2} \right)\).\\ We denote it by \(
    \varphi(X_1, X_2) \).
\end{definition}
\begin{remark}
    This definition follows naturally from the Euclidean law of cosines,
    namely \( c^{2} = a^2 + b^2 - 2ab\cos (a,b) \).
\end{remark}
Before formulating the property of the angle, we need the following auxiliary
lemma.
\begin{lemma}
    \label{lem:lambdacloud} For any cloud \( [M] \) and \( \lambda
    \in\mathbb{R}^+ \), the following properties hold:
    \begin{enumerate}
        \item \( \St\big(\lambda [M]\big) = \St\big([M]\big) \),
            \label{lem:lambdacloud:1}
        \item \( \lambda [M] = [\lambda M] \).\label{lem:lambdacloud:2}
    \end{enumerate}
\end{lemma}
\begin{proof}
    (\ref{lem:lambdacloud:1}) Let \( \mu \in \St\big([M]\big) \). Then \( \mu
    \cdot \lambda [M] = \lambda \cdot \mu [M]= \lambda [M] \), i.e. \( \mu \in
    \St\big(\lambda [M]\big) \). Hence \( \St\big([M]\big) \subseteq
    \St\big(\lambda [M]\big) \). The reverse inclusion is obtained by
    considering
    the cloud \( [M] \) as \(\frac{1}{\lambda}\cdot \lambda [M] \).

    (\ref{lem:lambdacloud:2}) Since the stabilizers coincide, for any
    \( \mu \in  \St\big(\lambda [M])\) we have \( \mu \cdot \lambda M =
    \lambda \cdot \mu M = \lambda M \). Hence, \( \lambda M \) is the center
    of the cloud \( \lambda [M] \).
\end{proof}
The angle between spaces has the following property.
\begin{lemma}
    \label{lemmaAngleBetweenSpaces} For any \( X_1, X_2 \in [M] \), \( \lambda
    \in\mathbb{R}^+ \), we have \( \varphi (X_1, X_2) = \varphi (\lambda
            X_1,
    \lambda X_2) \).
\end{lemma}
\begin{proof}
    By Lemma \ref{lem:lambdacloud}, the space \( \lambda M \) is the
    center of \( \lambda
        [M]
    \). As in the definition of the angle, set \[ |X_1 M| = r_1,\quad
        | X_2 M | =r_2,\quad
        |X_1
    X_2| = d.\] By the definition of the angle, we obtain the following chain
    of equalities: \[ \varphi (\lambda X_1, \lambda X_2) =
        \arccos\left(\frac{\lambda
            ^{2} r_1^2 + \lambda ^2 r_2^2 - \lambda^2 d^2}{2 \lambda  r_1
        \lambda r_2}\right) =
        \varphi (X_1, X_2).
    \]
    This completes the proof.
\end{proof}
We now consider the two definitions of the angle of a cloud that interest us.
\begin{definition}
    The \emph{angle} of a cloud \( [M] \) is the quantity
    \[
        \varphi \big([M]\big)= \sup \big\{\varphi (X_1, X_2) \colon | X_1,M
        |, |X_2,M| \neq 0\big\}.
    \]
\end{definition}
\begin{definition}
    The \emph{isosceles angle} of a cloud \( [M] \) is the quantity
    \[
        \varphi_e \big([M]\big) = \sup \big\{\varphi (X_1, X_2) \colon |X_1,M |
        = |X_2,M| \neq 0\big\}.
    \]
\end{definition}
For these definitions, we obtain the following corollary from Lemma
\ref{lemmaAngleBetweenSpaces}.
\begin{lemma}
    For any cloud \( [M] \) and \( \lambda \in\mathbb{R}^+ \), we
    have:
    \label{lem:angles}
    \begin{enumerate}
        \item \(\varphi \big(\lambda [M]\big) = \varphi
            \big([M]\big)\),\label{lem:angles:1}
        \item \(\varphi_e \big(\lambda [M]\big) = \varphi_e
            \big([M]\big)\),\label{lem:angles:2}
        \item \( \varphi_e \big([M]\big) \le \varphi \big([M]\big)
            \),\label{lem:angles:3}
        \item \(0 \le \varphi ([M]) \le \pi\),\label{lem:angles:4}
        \item \(0 \le \varphi_e ([M]) \le \pi\).\label{lem:angles:5}
    \end{enumerate}
\end{lemma}

We give examples of angles of clouds.
\begin{lemma}
    For the clouds \( [\Delta _{1}], [\mathbb{R}] \), the following
    equalities hold: \label{lem:angleexam}
    \begin{enumerate}
        \item \( \varphi \big([\Delta _{1}]\big) = \frac{\pi }{2}
            \),\label{lem:angleexam:1}
        \item \( \varphi_e \big([\Delta _{1}]\big) = \frac{\pi }{3}
            \),\label{lem:angleexam:2}
        \item \( \varphi \big([\mathbb{R}]\big) = \varphi_e
            \big([\mathbb{R}]\big) = \pi  \).\label{lem:angleexam:3}
    \end{enumerate}
\end{lemma}
\begin{proof}
    (\ref{lem:angleexam:1}) From the ultrametric inequality, it follows that \(
    \varphi \big([\Delta _1]\big)\le \frac{\pi }{2} \). To prove the reverse
    inequality, we find the angle between the two-point simplex \( \Delta
    _{2} = \{x_1, x_2\} \) and the three-point simplex \(
    \lambda \Delta _3  = \{y_1,y_2, y_3\}\). We use the formula proved in
    \cite{ivanov2016} for the Gromov--Hausdorff distance
    between two three-point spaces:
    \[
        d_{GH}(X,Y) = \frac{1}{2} \max\big\{ |a_1 - a_2|,|b_1 - b_2|,|c_1
        - c_2|\big\}.
    \]
    Here \( a_1 \le b_1 \le c_1 \) are the distances between points in
    the space \( X \), and \( a_2 \le b_2 \le c_2 \) are the distances between
    points in the space \( Y \). In our case, these are \( 0,1,1 \) and \(
    \lambda , \lambda , \lambda . \) Accordingly,
    \[
        d_{GH}(\Delta_2, \lambda \Delta _{3}) =\frac{1}{2}
        \max\big\{|\lambda |, | \lambda  -1|, |\lambda -1|\big\},
    \]
    which
    for \( \lambda \ge \frac 1 2 \) equals \(\frac \lambda 2 \). Thus, \(d
        _{GH}(\Delta _2, \lambda \Delta _3) =
    \frac{\lambda }{2} \). Next, taking into account that \( |\Delta _1
        \Delta _2|=\frac{1}{2}, |\Delta _1 \lambda \Delta
    _3|=\frac{\lambda }{2} \), we find the angle between \( \Delta _2,
    \lambda \Delta _3 \):
    \[ \cos \big(\varphi (\Delta _2 , \lambda \Delta _3)\big)
        =\frac{\frac{\lambda ^2}{4} + \frac{1}{4} - \frac{\lambda
        ^2}{4}}{\frac{\lambda }{2}} = \frac{1}{2 \lambda }
        \rightarrow 0,\quad \lambda \rightarrow \infty.
    \]
    Hence, \( \varphi \big([\Delta _1]\big) \ge \sup \varphi (\Delta _2,
    \lambda \Delta _3) = \frac \pi 2. \)

    (\ref{lem:angleexam:2}) From the ultrametric inequality, it follows that \(
    \varphi_e \big([\Delta _1]\big)\le \frac{\pi }{3} \). To prove the reverse
    inequality, it suffices to find spaces of equal diameter whose distance is
    equal to their diameters. Such spaces are, for example, \( \Delta
    _2, \Delta _3 \).

    (\ref{lem:angleexam:3}) It suffices to give examples of spaces \( X,Y\in
    [\mathbb{R}] \) such that \( |X \mathbb{R}| = |Y \mathbb{R}| = \frac 1 2
    |XY|. \) In \cite{mikhailovUnreleased}, it was proved that such
    spaces are \( \mathbb{Z}, \mathbb{R}\times
    [0,1] \).
    %Insert reference to Ivan with title.

\end{proof}
